\section{Kết luận chung}

Đồ án đã thực hiện quá trình phân tích doanh số và xu hướng tiêu dùng khách hàng trong lĩnh vực bán lẻ dựa trên bộ dữ liệu Superstore Sales. Thông qua các bước tiền xử lý dữ liệu, phân tích khám phá, trực quan hóa bằng Power BI và xây dựng các mô hình học máy, đề tài đã làm rõ được một số đặc điểm quan trọng trong hoạt động kinh doanh bán lẻ.

Kết quả phân tích cho thấy doanh thu, lợi nhuận và hành vi tiêu dùng của khách hàng có sự khác biệt theo thời gian, khu vực, danh mục sản phẩm và phân khúc khách hàng. Một số nhóm sản phẩm và khu vực có khả năng đóng góp lợi nhuận tốt, trong khi một số nhóm khác có doanh thu cao nhưng hiệu quả lợi nhuận chưa tương xứng. Đặc biệt, yếu tố chiết khấu có ảnh hưởng đáng kể đến lợi nhuận, cho thấy doanh nghiệp cần kiểm soát chính sách chiết khấu một cách hợp lý.

Bên cạnh đó, kết quả xây dựng mô hình cho thấy các mô hình học máy phi tuyến như Random Forest, Gradient Boosting và XGBoost cho hiệu quả dự báo tốt hơn so với hồi quy tuyến tính đa biến. Trong đó, XGBoost là mô hình đạt kết quả tốt nhất trên tập kiểm tra. Ngoài ra, mô hình K-Means Clustering cũng hỗ trợ phân nhóm khách hàng theo doanh thu và lợi nhuận, góp phần bổ sung góc nhìn về hành vi và giá trị khách hàng.

Nhìn chung, đề tài đã đạt được mục tiêu đề ra là khai thác dữ liệu bán hàng để phân tích tình hình kinh doanh, nhận diện xu hướng tiêu dùng và ứng dụng các mô hình phân tích dữ liệu nhằm hỗ trợ đánh giá, dự báo và ra quyết định dựa trên dữ liệu.

\section{Những kết quả đạt được}

Trong quá trình thực hiện, đồ án đã đạt được một số kết quả chính như sau:

\begin{itemize}
	\item Hệ thống hóa được cơ sở lý thuyết liên quan đến phân tích dữ liệu trong kinh doanh, trực quan hóa dữ liệu, tiền xử lý dữ liệu và các mô hình học máy sử dụng trong đề tài.
	
	\item Thực hiện quá trình tiền xử lý dữ liệu, kiểm tra dữ liệu thiếu, dữ liệu trùng lặp, chuẩn hóa kiểu dữ liệu và xử lý ngoại lai, giúp dữ liệu phù hợp hơn cho phân tích và xây dựng mô hình.
	
	\item Phân tích khám phá dữ liệu nhằm làm rõ đặc điểm doanh thu, lợi nhuận, số lượng bán, chiết khấu, khu vực bán hàng, danh mục sản phẩm và phân khúc khách hàng.
	
	\item Xây dựng dashboard trên Power BI để trực quan hóa các chỉ tiêu quan trọng, giúp quan sát tình hình kinh doanh một cách rõ ràng và trực quan hơn.
	
	\item Xây dựng và đánh giá các mô hình dự báo lợi nhuận gồm hồi quy tuyến tính đa biến, Random Forest, Gradient Boosting và XGBoost. Kết quả cho thấy XGBoost là mô hình phù hợp nhất trong phạm vi đề tài.
	
	\item Ứng dụng mô hình K-Means Clustering để phân cụm khách hàng, từ đó hỗ trợ nhận diện các nhóm khách hàng có đặc điểm khác nhau về doanh thu và lợi nhuận.
	
	\item Rút ra được một số nhận xét có ý nghĩa thực tiễn, đặc biệt là vai trò của chiết khấu, doanh thu, số lượng bán và nhóm sản phẩm đối với lợi nhuận trong dữ liệu bán lẻ.
\end{itemize}

\section{Hạn chế tồn đọng của đề tài}

Mặc dù đã đạt được một số kết quả nhất định, đề tài vẫn còn tồn tại một số hạn chế. Trước hết, bộ dữ liệu sử dụng là dữ liệu có sẵn nên chưa phản ánh đầy đủ toàn bộ các yếu tố thực tế có thể ảnh hưởng đến hoạt động kinh doanh, chẳng hạn như chi phí vận hành, chính sách marketing, hành vi khách hàng chi tiết hoặc yếu tố cạnh tranh thị trường.

Thứ hai, các mô hình dự báo trong đề tài mới được xây dựng ở mức cơ bản, chủ yếu nhằm so sánh hiệu quả giữa các thuật toán. Việc tối ưu tham số mô hình chưa được thực hiện chuyên sâu, do đó kết quả dự báo vẫn có thể được cải thiện thêm nếu áp dụng các kỹ thuật tối ưu phù hợp.

Thứ ba, mô hình K-Means mới chỉ phân cụm khách hàng dựa trên một số đặc điểm chính như doanh thu và lợi nhuận. Trong thực tế, việc phân nhóm khách hàng có thể cần kết hợp thêm nhiều yếu tố khác như tần suất mua hàng, giá trị trung bình mỗi đơn hàng, thời gian mua gần nhất và hành vi mua theo từng nhóm sản phẩm.

Ngoài ra, dashboard trong đề tài chủ yếu phục vụ mục tiêu trực quan hóa và phân tích mô tả. Việc tích hợp trực tiếp kết quả dự báo từ mô hình vào dashboard chưa được triển khai, nên khả năng hỗ trợ ra quyết định theo thời gian thực vẫn còn hạn chế.

\section{Hướng phát triển trong tương lai}

Trong tương lai, đề tài có thể được phát triển theo một số hướng sau:

\begin{itemize}
	\item \textbf{Mở rộng dữ liệu nghiên cứu:} 
	Có thể bổ sung thêm các dữ liệu liên quan đến chi phí vận hành, chương trình khuyến mãi, hành vi khách hàng và dữ liệu thị trường. Việc mở rộng dữ liệu giúp kết quả phân tích phản ánh sát hơn bối cảnh kinh doanh thực tế.
	
	\item \textbf{Tối ưu mô hình dự báo:} 
	Các mô hình học máy có thể được cải thiện bằng cách tối ưu tham số thông qua Grid Search, Random Search hoặc Cross Validation. Điều này giúp nâng cao độ chính xác dự báo và khả năng tổng quát hóa trên dữ liệu mới.
	
	\item \textbf{Phát triển phân cụm khách hàng chuyên sâu hơn:} 
	Đối với bài toán phân cụm, có thể kết hợp thêm phương pháp RFM gồm Recency, Frequency và Monetary Value. Cách tiếp cận này giúp phân nhóm khách hàng chi tiết hơn, hỗ trợ tốt hơn cho hoạt động chăm sóc khách hàng và marketing.
	
	\item \textbf{Nâng cấp dashboard trực quan:} 
	Dashboard Power BI có thể được phát triển thêm theo hướng tương tác sâu hơn, cho phép người dùng theo dõi dữ liệu theo thời gian, khu vực, nhóm sản phẩm và phân khúc khách hàng. Ngoài ra, có thể tích hợp kết quả dự báo từ mô hình vào dashboard.
	
	\item \textbf{Xây dựng hệ thống phân tích dữ liệu tự động:} 
	Một hướng phát triển quan trọng là xây dựng thành website, phần mềm hoặc hệ thống hỗ trợ phân tích dữ liệu bán hàng tự động. Khi đó, người dùng chỉ cần tải bộ dữ liệu lên hệ thống, phần mềm sẽ tự động kiểm tra dữ liệu, làm sạch dữ liệu, trực quan hóa, xây dựng mô hình dự báo và xuất báo cáo phân tích cơ bản.
\end{itemize}

Tóm lại, đề tài là bước đầu trong việc vận dụng phân tích dữ liệu, trực quan hóa và học máy vào bài toán kinh doanh bán lẻ. Các kết quả đạt được có thể làm nền tảng để tiếp tục phát triển các hệ thống phân tích, dự báo và hỗ trợ ra quyết định trong thực tiễn.	